% ---------------- 附录一：器物与编号对照 ----------------
\chapter*{附录一\quad 器物与编号对照}
\phantomsection
\addcontentsline{toc}{chapter}{附录一\quad 器物与编号对照}
\markboth{附录一}{附录一}
\begingroup
\footnotesize
\linespread{1.2}\selectfont
\begin{longtable}{@{}p{30mm}p{22mm}p{16mm}p{26mm}p{28mm}@{}}
\toprule
简次 & 编号 & 长 (cm) & 材质 & 备注\\
\midrule
\endfirsthead
\toprule
简次 & 编号 & 长 (cm) & 材质 & 备注\\
\midrule
\endhead
简一 & 11 & 21.3 & 金属 & \\
简二 & 13 & 17.6 & 陶瓷 & 有釉\\
简三 & 15 & 24.1 & 金属 & \\
简四 & 18 & 19.8 & 金属 & 右端断裂\\
简五 & 21 & 22.4 & 陶瓷 & \\
简六 & 24 & 15.2 & 金属 & 下端蚀\\
简七 & 28 & 11.7 & 陶瓷 & 出土地点自记\\
简八 & 31 & 18.9 & 金属 & \\
简九 & 34 & 26.8 & 金属 & 卷一最长\\
简十 & 37 & 20.5 & 金属 & \\
简十一 & 39 & 14.8 & 陶瓷 & \\
简十二 & 41 & 9.4 & 金属 & 全书最短\\
简十三 & 43 & 17.1 & 金属 & \\
简十四 & 26 & 18.4 & 金属 & 二次编目，次序错乱\\
简十五 & 45 & 13.6 & 金属 & \\
简十六 & 47 & 16.9 & 陶瓷 & \\
简十七 & 49 & 12.3 & 金属 & \\
简十八 & 51 & 15.8 & 金属 & \\
简十九 & 53 & 19.2 & 金属 & 右端灼痕\\
简二十 & 57 & 22.0 & 金属 & \\
简二十一 & 60 & 28.4 & 金属 & 全书最长\\
简二十二 & 61 & 23.7 & 金属 & 字浅\\
简二十三 & 63 & 16.4 & 陶瓷 & \\
简二十四 & 65 & 20.1 & 金属 & \\
简二十五 & 67 & 14.2 & 陶瓷 & \\
简二十六 & 69 & 18.7 & 金属 & 铅笔痕\\
简二十七 & 71 & 15.3 & 金属 & \\
简二十八 & 73 & 10.8 & 金属 & 下端蚀\\
简二十九 & 75 & 25.6 & 金属 & \\
简三十 & 77 & 17.9 & 陶瓷 & \\
简三十一 & 79 & 12.9 & 金属 & \\
简三十二 & 81 & 13.4 & 金属 & 字浅，后世题\\
简三十三 & 83 & 18.1 & 金属 & 疑摹刻\\
简三十四 & 85 & 11.2 & 陶瓷 & 背面朱印\\
简三十五 & 87 & 9.8 & 金属 & 上下俱蚀\\
简三十六 & 89 & 12.6 & 金属 & 所述圆片存 ФВ-1963-46\\
简三十七 & 91 & 6.2 & 金属 & 末段\\
\bottomrule
\end{longtable}
\endgroup

\vspace{3mm}
\noindent{\footnotesize\color{pencil}
另：存疑八条所用原件分别为 12、29、36、44、55、62、68、88 号，不在上表所列三十七片之内，故三十七片与八条合计四十五件，而全书残片共八十九号。

\chapter*{附录二\quad 年表}
\phantomsection
\addcontentsline{toc}{chapter}{附录二\quad 年表}
\markboth{附录二}{附录二}
\begingroup
\footnotesize
\linespread{1.2}\selectfont
\begin{longtable}{@{}p{30mm}p{72mm}p{32mm}@{}}
\toprule
年代 & 事件 & 依据\\
\midrule
\endfirsthead
\toprule
年代 & 事件 & 依据\\
\midrule
\endhead
前 421 年 & 第一批撤离（约 6\,000 人） & 简二十九、简三十\\
前 380 年 & 第二批（约 8\,000 人） & 简二十九\\
前 340 年 & 第三批（约 10\,000 人） & 简二十九\\
前 300 年 & 第四批（约 13\,000 人） & 简二十九\\
前 260 年 & 第五批（约 18\,000 人） & 简二十九\\
前 221 年 & 第六批（65\,000 人）；城封存 & 简二十、简二十一至二十四\\
前 221 年 & 地热衰竭，气路水路断绝 & 简二十七\\
公元 151 年 & 「后三百七十二年，有客至」 & 简三十二（存疑）\\
一五八九年 & 内官监采办「至此而返」 & 简三十四（存疑）\\
一九六三年 & 残片三十七件编目 & ФВ-1963-11 至 47\\
\bottomrule
\end{longtable}
\endgroup

\chapter*{附录三\quad 引用书目}
\phantomsection
\addcontentsline{toc}{chapter}{附录三\quad 引用书目}
\markboth{附录三}{附录三}
\begingroup
\footnotesize
\linespread{1.25}\selectfont
\begin{list}{}{\leftmargin=1.8em \itemindent=-1.8em \topsep=0pt \parsep=1.4pt \itemsep=0pt}
\item 《史记》，中华书局 1959 年点校本。
\item 《汉书》，中华书局 1962 年点校本。
\item 《左传》，杨伯峻注本，中华书局 1990 年版。
\item 《国语》，上海古籍出版社 1978 年点校本。
\item 《论衡》，中华书局 1990 年点校本。
\item 《古事记》，人民文学出版社 1979 年版。
\item 《日本书纪》，人民出版社「东方文化集成」本。
\item Баранников А. В. Лунная прародина: опыт реконструкции. М.: Изд-во МГУ, 2018.
\item Баранников А. В. Восстановление антиматериальной сети Чанъани. М.: Изд-во МГУ, 2019.
\item Баранников А. В. Атлас реконструкций лунных объектов. СПб., 2016.
\item Баранников А. В. 战争之雷霆（俄文原著 2019 年；中译 2021 年据全文重译；第四版 2024 年）。
\end{list}
\endgroup

% ---------------- 索引 ----------------
\clearpage
\chapter*{索\quad 引}
\phantomsection
\addcontentsline{toc}{chapter}{索\quad 引}
\markboth{索\quad 引}{索\quad 引}
\begingroup
\footnotesize
\linespread{1.35}\selectfont
\newcommand{\idx}[2]{\hangindent=1.4em \noindent #1\dotfill #2\par}
\idx{六批}{14、15、30、31、38}
\idx{净秽}{5、7、8、24、25、27}
\idx{净秽山}{5、8、25、27}
\idx{圆片}{5、17、37}
\idx{地热之衰}{14、32}
\idx{外工}{7、10、14、21、23、29、32}
\idx{存活率}{14、22、33}
\idx{小篆}{1}
\idx{居住环}{1、5、6、7、12}
\idx{岁计}{5、9、11、14、15、27、29、31}
\idx{广寒宫}{1、16}
\idx{月之民}{21、22、33}
\idx{月兔}{10、11、25、27}
\idx{玉山}{17}
\idx{界}{3、5、6、8、12、13、24、27}
\idx{界石}{13}
\idx{纯化}{25}
\idx{蓬莱药}{25}
\idx{退役清单}{2、12}
\idx{静海}{1、3}
\endgroup
